% !TEX TS-program = pdflatex
% !TEX encoding = UTF-8 Unicode

% This file is a template using the "beamer" package to create slides for a talk or presentation
% - Giving a talk on some subject.
% - The talk is between 15min and 45min long.
% - Style is ornate.

% MODIFIED by Jonathan Kew, 2008-07-06
% The header comments and encoding in this file were modified for inclusion with TeXworks.
% The content is otherwise unchanged from the original distributed with the beamer package.

\documentclass{beamer}
\setbeamercovered{invisible}



% Copyright 2004 by Till Tantau <tantau@users.sourceforge.net>.
%
% In principle, this file can be redistributed and/or modified under
% the terms of the GNU Public License, version 2.
%
% However, this file is supposed to be a template to be modified
% for your own needs. For this reason, if you use this file as a
% template and not specifically distribute it as part of a another
% package/program, I grant the extra permission to freely copy and
% modify this file as you see fit and even to delete this copyright
% notice. 


\mode<presentation>
{
  \usetheme{Warsaw}
  % or ...

  % or whatever (possibly just delete it)
}


\usepackage[english]{babel}
% or whatever

\usepackage[utf8]{inputenc}
% or whatever

\usepackage{times}
\usepackage[T1]{fontenc}
% Or whatever. Note that the encoding and the font should match. If T1
% does not look nice, try deleting the line with the fontenc.

%%% MATH RELATED

\input{macros.tex}

\title{MATH 350-2 Advanced Calculus} 
\subtitle
{} % (optional)

\author[W.R. Casper] % (optional, use only with lots of authors)
{W.R. Casper}
% - Use the \inst{?} command only if the authors have different
%   affiliation.

\institute[California State University Fullerton] % (optional, but mostly needed)
{
  Department of Mathematics\\
  California State University Fullerton}
% - Use the \inst command only if there are several affiliations.
% - Keep it simple, no one is interested in your street address.

\subject{Talks}
% This is only inserted into the PDF information catalog. Can be left
% out. 



% If you have a file called "university-logo-filename.xxx", where xxx
% is a graphic format that can be processed by latex or pdflatex,
% resp., then you can add a logo as follows:

% \pgfdeclareimage[height=0.5cm]{university-logo}{university-logo-filename}
% \logo{\pgfuseimage{university-logo}}



% Delete this, if you do not want the table of contents to pop up at
% the beginning of each subsection:
\AtBeginSubsection[]
{
  \begin{frame}<beamer>{Outline}
    \tableofcontents[currentsection,currentsubsection]
  \end{frame}
}


% If you wish to uncover everything in a step-wise fashion, uncomment
% the following command: 

%\beamerdefaultoverlayspecification{<+->}


\begin{document}

\begin{frame}
  \titlepage
\end{frame}

\begin{frame}{Outline}
  \tableofcontents
  % You might wish to add the option [pausesections]
\end{frame}

% Since this a solution template for a generic talk, very little can
% be said about how it should be structured. However, the talk length
% of between 15min and 45min and the theme suggest that you stick to
% the following rules:  

% - Exactly two or three sections (other than the summary).
% - At *most* three subsections per section.
% - Talk about 30s to 2min per frame. So there should be between about
%   15 and 30 frames, all told.

\section{Real Analysis Lecture 3}

\subsection{Set basics}

\begin{frame}{Set basics}
Intuitively, a \textbf{set} $A$ is a "collection of things" which we call the \textbf{elements} of $A$.
\begin{itemize}
\pause
\item in practice, this is a bad definition (Russell's Paradox)
\pause
\item true set formulation: Zermelo-Frankel Axioms
\end{itemize}
\pause
Examples:
\begin{itemize}
\pause
\item $\mathbb R$, $\mathbb{Z}_+$, $\mathbb{Z}$, $\mathbb{Q}$
\pause
\item $(1,5]$, $(0,\infty)$
\pause
\item empty set $\varnothing$
\pause
\item $\{\heartsuit,\text{Fall}, \{\varnothing\}\}$
\pause
\item $\{n\in \mathbb{Z}: \text{$n$ is prime}\}$
\end{itemize}
\end{frame}

\begin{frame}{Everything is a set}
In the true minimalistic philosophy of mathematics, we want everything to be a set.
\begin{itemize}
\pause
\item integers
\pause
$$0 = \varnothing, 1 = \{\varnothing\}, 2 = \{\varnothing,\{\varnothing\}\},\dots$$
\pause
\item ordered pairs
\pause
$$(a,b) = \{a,\{a,b\}\}$$
\pause
\item even relations and functions are sets!
\end{itemize}
\end{frame}

\begin{frame}{Challenge!}
\begin{prob}
Prove that for ordered pairs $(a,b)$ and $(c,d)$ that
$$(a,b)=(c,d)\quad\text{if and only if}\quad a=c\ \text{and}\ b=d$$
\end{prob}
\pause
\begin{soln}
$(a,b)=(c,d)$ if and only if $\{a,\{a,b\}\}=\{c,\{c,d\}\}$\\
\pause
Clearly, if $a=c$ and $b=d$, then $\{a,\{a,b\}\}=\{c,\{c,d\}\}$\\
\pause
The tough part is the opposite direction!\\
\pause
Suppose $\{a,\{a,b\}\}=\{c,\{c,d\}\}$.\\
\pause
Two possible cases:
\pause
\begin{align*}
\text{Case I:}\quad a=c\ \ \text{and}\ \  \{a,b\} = \{c,d\}\\
\pause
\text{Case II:}\quad a=\{c,d\}\ \ \text{and}\ \ \{a,b\} = c
\end{align*}
\end{soln}
\end{frame}

\begin{frame}{Challenge!}
\begin{prob}
Prove that for ordered pairs $(a,b)$ and $(c,d)$ that
$$(a,b)=(c,d)\quad\text{if and only if}\quad a=c\ \text{and}\ b=d$$
\end{prob}
\begin{soln}
$$\text{Case I:}\quad a=c\ \ \text{and}\ \  \{a,b\} = \{c,d\}$$
\pause
Since $\{a,b\} = \{c,d\}$, we know $b\in \{c,d\}$.\\
\pause
Therefore $b=c$ or $b=d$.\\
If $b=d$, we're done! \pause ... so assume instead that $b=c$\\
\pause
Then $a=c$ implies $a=b$.\\
\pause
Therefore $\{c,d\} = \{a,b\} = \{a,a\} = \{a\}$.\\
\pause
It follows that $d=c=b=a$.
\end{soln}
\end{frame}

\begin{frame}{Challenge!}
\begin{prob}
Prove that for ordered pairs $(a,b)$ and $(c,d)$ that
$$(a,b)=(c,d)\quad\text{if and only if}\quad a=c\ \text{and}\ b=d$$
\end{prob}
\begin{soln}
$$\text{Case II:}\quad a=\{c,d\}\ \ \text{and}\ \ \{a,b\} = c$$
\pause
This would imply that $c\in a$ and $a\in c$.\\
\pause
This can be shown to contradict the ZF Axioms of Set Theory.\\
\pause
Specifically the regularity axiom for the set $\{a,c\}$...
\end{soln}
\end{frame}

\subsection{Set algebra}

\begin{frame}{Unions of sets}
Unions:
\pause
$$A\cup B = \{x: x\in A\ \ \text{or}\ \ x\in B\},$$
\pause
$$A\cup B\cup C = \{x: x\in A\ \ \text{or}\ \ x\in B\ \ \text{or}\ \ x\in C\},$$
\pause
$$A_1\cup A_2\cup \dots \cup A_n = \{x: x\in A_i\ \ \text{for some $1\leq i\leq n$}\}$$
\pause
{\color{red}NOTATION:}
$$\bigcup_{i=1}^n A_i = A_1\cup A_2\cup \dots \cup A_n.$$
\end{frame}

\begin{frame}{Intersections of sets}
Intersections:
\pause
$$A\cap B = \{x: x\in A\ \ \text{and}\ \ x\in B\},$$
\pause
$$A\cap B\cap C = \{x: x\in A\ \ \text{and}\ \ x\in B\ \ \text{and}\ \ x\in C\},$$
\pause
$$A_1\cap A_2\cap \dots \cap A_n = \{x: x\in A_i\ \ \text{for all $1\leq i\leq n$}\}$$
\pause
{\color{red}NOTATION:}
$$\bigcap_{i=1}^n A_i = A_1\cap A_2\cap \dots \cap A_n.$$
\end{frame}

\begin{frame}{Complements of sets}
\pause
The \vocab{complement} of $A$ relative to $B$ is
$$B-A = \{b\in B: b\notin A\}.$$
\pause
\begin{thm}{De Morgan's Laws}
\pause
$$B-\bigcup_{i=1}^n A_i = \bigcap_{i=1}^n (B-A_i)$$
\pause
$$B-\bigcap_{i=1}^n A_i = \bigcup_{i=1}^n (B-A_i)$$
\end{thm}
\end{frame}


\begin{frame}{Families of sets}
A \textbf{family of sets} is a collection
$$\{A_i: i\in I\}$$
\begin{itemize}
\pause
\item $I$ any set, called the \textbf{index set}
\pause
\item $A_i$ is a set for all $i$
\end{itemize}
\pause
$$\bigcup_{i\in I} A_i = \{x: x\in A_i,\ \text{for some $i\in I$}\}$$
\pause
$$\bigcap_{i\in I} A_i = \{x: x\in A_i,\ \text{for all $i\in I$}\}$$
\end{frame}

\begin{frame}{Challenge}
\begin{itemize}
\item index set $I = (1,2)$
\item family of sets $\{A_i: i\in I\}$
\item $A_i = [0,i]$
\end{itemize}
\begin{prob}
Determine $\bigcup_{i\in I} A_i$.
\end{prob}
\end{frame}

\begin{frame}{Challenge}
\begin{itemize}
\item index set $I = \mathbb{Z}_+$
\item family of sets $\{A_i: i\in I\}$
\item $A_i = [0,i)$
\end{itemize}
\begin{prob}
Determine $\bigcap_{i\in I} A_i$.
\end{prob}
\end{frame}

\begin{frame}{Challenge}
\begin{prob}
Prove that De Morgan's Laws apply to families of sets too, ie.
$$B-\bigcup_{i\in I} A_i = \bigcap_{i\in I} (B-A_i)$$
\pause
$$B-\bigcap_{i\in I} A_i = \bigcup_{i\in I} (B-A_i)$$
\end{prob}
\end{frame}

\begin{frame}{Relations}
The \textbf{Cartesian product} of $A$ and $B$ is
$$A\times B = \{(a,b): a\in A,\ b\in B\}.$$
\pause
A \textbf{relation} $\mathcal R$ from $A$ to $B$ is a subset of $A\times B$.\\
\pause
$$\text{\color{red}NOTATION:}\quad a\mathcal R b\quad\text{means}\quad (a,b)\in \mathcal R.$$
\pause
Domain and codomain:
\pause
\begin{align*}
\pause
\text{dom}(\mathcal{R}) &= \{a\in A: \exists b\in B,\ a\mathcal R b\}\\
\text{codom}(\mathcal{R}) &= \{b\in B: \exists a\in A,\ a\mathcal R b\}
\end{align*}
\end{frame}

\begin{frame}{Relations}
A relation $\mathcal R$ from $A$ to $A$ is called a \textbf{relation on} $A$.  A relation on $A$ is
\begin{itemize}
\pause
\item \textbf{reflexive} if $a\mathcal Ra$ for all $a\in A$
\pause
\item \textbf{symmetric} if $a\mathcal Rb$ implies $b\mathcal Ra$ for all $a,b\in A$
\pause
\item \textbf{transitive} if $a\mathcal Rb$ and $b\mathcal Rc$ implies $a\mathcal Rc$ for all $a,b,c\in\mathcal A$
\end{itemize}
\pause
An \textbf{equivalence relation} satisfies all three properties.\\
\pause
Examples:
\pause
\begin{itemize}
\pause
\item $<\ = \{(x,y): y-x\in (0,\infty)\}$ is transitive but not reflexive or symmetric on $\mathbb{R}$
\pause
\item $\leq\ = \{(x,y): y-x\in [0,\infty)\}$ is reflexive and transitive but not symmetric on $\mathbb{R}$
\end{itemize}
\end{frame}

\begin{frame}{Challenge!}
\begin{prob}
Give an example of a relation on $\mathbb{R}$ which is symmetric and transitive but not reflexive.
\end{prob}
\end{frame}

\begin{frame}{Challenge!}
\begin{prob}
Give an example of a relation on $\mathbb{R}$ which is reflexive and symmetric but not transitive.
\end{prob}
\end{frame}

\begin{frame}{Functions}
A \textbf{function} from $A$ to $B$ is a relation $\mathcal R$ from $A$ to $B$ with the property
\pause
$$a\mathcal Rb\ \text{and}\ a\mathcal Rc\quad\text{implies}\quad b=c.$$
\pause
If a relation $\mathcal R$ is a function, we usually use a symbol like $f$.\\
\pause
\begin{align*}
\text{\color{red}NOTATION:}&\quad f: A\rightarrow B\quad\text{means $f$ is a function from $A$ to $B$}\\
\text{\color{red}NOTATION:}&\quad f(a) = b\quad\text{means}\quad (a,b)\in f.
\end{align*}
\pause
The set $$\img(f) = \{f(a): a\in A\}$$ is called the \textbf{range} or \textbf{image} of $f$
\end{frame}

\begin{frame}{Challenge!}
\begin{prob}
Determine all the equivalence relations on $\mathbb{R}$ which are also functions.
\end{prob}
\pause
\begin{soln}
$f=\mathcal R$ must be reflexive, so $x\mathcal Rx$ for all $x$\\
\pause
This means $f(x) = x$ for all $x$.\\
\pause
Thus the only function which is an equivalence relation is the identity function
$$f(x) = x.$$
\end{soln}
\end{frame}

\end{document}


